\section{性能与时序优化}\label{sec:optimization}

\subsection{优化方法}

优化以完整工作负载、周期级 RTL 统计和 Vivado 布局布线报告为依据。每项成果必须同时满足功能语义
和明确的数据口径：性能比较只在同一镜像内部进行；时序数字注明频率及 routed/post-route physopt
阶段；不同提交或不同 COE 的数据不拼接为一个候选结果。本文只描述冻结版本仍保留的成功结构。

\begin{figure}[H]\centering
\begin{tikzpicture}[s/.style={draw,rounded corners=1.2pt,minimum width=2.55cm,minimum height=0.78cm,align=center,fill=blue!6},a/.style={-Stealth,thick},font=\small]
\node[s] (measure) {完整运行与计数器};\node[s,right=0.75cm of measure] (locate) {定位高频停顿/关键路径};
\node[s,right=0.75cm of locate] (rtl) {结构化 RTL 调整};\node[s,below=0.8cm of rtl] (verify) {功能与周期复核};
\node[s,left=0.75cm of verify] (impl) {Vivado 实现与 WNS};\node[s,below=0.8cm of measure] (record) {证据归档与结果复核};
\draw[a](measure)--(locate);\draw[a](locate)--(rtl);\draw[a](rtl)--(verify);\draw[a](verify)--(impl);\draw[a](impl)--(record);\draw[a](record)--(measure);
\end{tikzpicture}\caption{以测量和物理结果驱动的优化闭环}\end{figure}


\subsection{前端控制流优化}

冻结版本将 BTB 扩大至 512 项，并使用八 bank LUTRAM 控制查询深度；更新端保留窄 shadow，避免更新
数据绕回取指查询通路。条件分支采用表项局部 2-bit 饱和计数器，JAL 和 return 使用无条件目标，return
还可由两项 RAS 提供不依赖 BTB 目标别名的地址。未见过且实际不跳转的条件分支不分配表项，从而保留
已有控制转移目标。重定向 fetch 固定预测顺序下一条，后续再恢复正常预测，避免重定向邮箱与异步查询
同拍串联。

\subsection{缓存与 load-use 优化}

数据缓存由 2\,KiB 扩展为 4\,KiB 双 bank，减少直接映射冲突。主数据继续使用 BRAM 以控制资源，
逐字节 LUTRAM 副本则提供异步窄加载数据。late-load 只开放给固定、容易验证且逻辑短的消费者；数据
在 EXU--LSU 或响应捕获边界注册后再参与分支比较和结果形成。store 能安全覆盖整字时直接分配缓存行，
窄 store 不能安全合并时保守失效；store epoch 防止旧 load 回填覆盖新数据。

\subsection{执行和后端通路优化}

普通整数 ALU、直接结果、长算术、领域加速和 load 结果分成五类。相邻快速结果采用分组编码选择，
不允许其控制信号进入地址、M/B、自定义指令和 CSR 等长路径。地址操作数单独注册，D-cache 查询只
携带所需 index/tag；加速单元的控制、响应和结果也分别经过边界寄存。该结构的目标不是增加流水深度，
而是让每类消费者只看到所需数据，降低宽 Mux、扇出和跨模块组合回路。

\subsection{乘法和位操作优化}

普通低位乘法使用较短 DSP 流水；操作数满足定宽条件时进入一周期窄乘法通路，高位乘法仍使用完整
有符号/无符号乘积。B 扩展按逻辑深度分为短组合操作和迭代操作，避免简单的移位加法或符号扩展承担
通用迭代延迟。每条多周期指令在结果未就绪期间持续提供目的寄存器占用信息，保持 RAW 停顿正确。

\subsection{分赛区后续性能演进}

表\ref{tab:regional-perf}用于吸收分赛区后续优化报告中仍有效的量化结果。两套工作负载彼此独立，
只能在各自表内纵向比较。数据没有用于推算总决赛冻结版本的板测成绩。

\begin{table}[H]\centering
\caption{分赛区后续版本的同口径性能演进}\label{tab:regional-perf}
\begin{tabularx}{\linewidth}{@{}Yp{0.19\linewidth}p{0.18\linewidth}p{0.20\linewidth}@{}}
\toprule
\textbf{同一工作负载内的成功里程碑}&\textbf{withMext-v2 时间}&\textbf{西北样例时间}&\textbf{说明}\\\midrule
分赛区主报告基线&2.020 s&---&280 MHz 统一复测\\
乘法与方向预测优化&1.943 s&---&减少乘法等待\\
D-cache 整字写分配&1.868 s&---&改善写后读局部性\\
load-use 晚数据旁路&1.643 s&---&受限消费者提前获得数据\\
多周期单元与地址通路&1.606910 s&6.866 s&时序与性能共同整理\\
窄加载与缓存 shadow&1.607404 s&6.553 s&不同程序的访存组成不同\\
32 项 BTB 与返回预测&1.607403 s&6.481 s&控制流容量和间接返回\\
部分写命中保留&1.607403 s&6.182 s&提升缓存局部性\\
分赛区后续冻结结果&1.607403 s&6.184 s&接口整理后的闭合版本\\
\bottomrule
\end{tabularx}
\end{table}

\begin{figure}[H]\centering
\begin{tikzpicture}\begin{axis}[width=0.92\linewidth,height=6.5cm,symbolic x coords={基线,乘法预测,DCache,晚旁路,多周期地址,Shadow,BTB返回,部分写,冻结},xtick=data,x tick label style={rotate=28,anchor=east,font=\scriptsize},ylabel={评估时间（s）},ymin=1.5,ymax=2.08,grid=major]
\addplot[blue!70!black,thick,mark=*] coordinates {(基线,2.020)(乘法预测,1.943)(DCache,1.868)(晚旁路,1.643)(多周期地址,1.606910)(Shadow,1.607404)(BTB返回,1.607403)(部分写,1.607403)(冻结,1.607403)};
\end{axis}\end{tikzpicture}\caption{withMext-v2 同口径性能走向}\end{figure}



\subsection{总决赛成功里程碑链}\label{sec:performance}

总决赛阶段的成功结构按能力链归纳为：分赛区 tag 基线、动态预测与更大 BTB、4\,KiB 双 bank
D-cache、窄乘法与 B 扩展分流、领域自定义指令、结果通路分区、后端消费者与控制边界注册，最终冻结
于 \commit{fabb599}。由于这些里程碑曾使用不同的程序映像、扩展组合和频率，本文不把历史运行时间
强行绘成同一折线。取得冻结提交的正式板测后，仅在下表填入同一提交的可比数据。

\begin{table}[H]\centering
\caption{总决赛冻结版本待补性能结果}
\begin{tabularx}{\linewidth}{@{}p{0.24\linewidth}p{0.16\linewidth}p{0.23\linewidth}Y@{}}
\toprule
\textbf{测试项}&\textbf{频率}&\textbf{结果}&\textbf{必须附带的证据}\\\midrule
正式性能工作负载&250 MHz&\MetricPending{运行时间}&UART 原始输出、CRC、bitstream SHA-256、COE SHA-256\\
正式性能工作负载&300 MHz&\MetricPending{运行时间}&UART 原始输出、CRC、bitstream SHA-256、COE SHA-256\\
重复上板稳定性&待定&\MetricPending{统计结果}&逐次 capture 日志与板卡编号、样本数、均值和极差\\
资源利用率&250/300 MHz&\MetricPending{资源数据}&对应实现的 LUT/FF/BRAM/DSP utilization report\\
\bottomrule
\end{tabularx}
\end{table}

\PlaceholderFigure{冻结版本性能与 WNS 综合图}{待获得同一提交的正式板测时间后，生成运行时间、频率、WNS 与资源利用率组合图。}

\subsection{实现频率与 WNS}\label{sec:timing}

冻结提交目前有三组可明确归属的实现结果。250\,MHz 正确 worktree 在 routed 阶段 WNS 为
$-0.087$\,ns，经 post-route physopt 后达到 $+0.009$\,ns，完成 setup 时序闭合\evidence{E-FAB-250}。
280\,MHz 的正确复测 routed WNS 为 $-0.664$\,ns，之后 Vivado 在 post-route physopt 阶段异常退出，
因此不存在该次运行的最终 physopt 数字\evidence{E-FAB-280}。300\,MHz 归档的 routed WNS 为
$-0.590$\,ns，post-route physopt 为 $-0.580$\,ns\evidence{E-FAB-300,E-FAB-300P}，表示高频实现能够
完成布局布线，但尚未满足 setup 时序。

\begin{table}[H]\centering
\caption{冻结提交在不同频率下的实现结果}
\begin{tabularx}{\linewidth}{@{}p{0.12\linewidth}p{0.17\linewidth}p{0.22\linewidth}Yp{0.16\linewidth}@{}}\toprule
\textbf{频率}&\textbf{Routed WNS}&\textbf{Post-physopt WNS}&\textbf{状态}&\textbf{证据}\\\midrule
250 MHz&$-0.087$ ns&$+0.009$ ns&时序闭合&E-FAB-250\\
280 MHz&$-0.664$ ns&不可用&仅保留 routed 结果&E-FAB-280\\
300 MHz&$-0.590$ ns&$-0.580$ ns&实现完成，setup 未闭合&E-FAB-300/P\\
\bottomrule\end{tabularx}
\end{table}

\begin{figure}[H]\centering
\begin{tikzpicture}\begin{axis}[ybar,width=0.82\linewidth,height=6.0cm,xtick={0,1,2},
xticklabels={{\shortstack{240 MHz\\最终 RTL}},{\shortstack{300 MHz\\前一基线}},{\shortstack{300 MHz\\最终 RTL}}},
ylabel={WNS（ns）},ymin=-0.82,ymax=0.12,bar width=22pt,
nodes near coords={\pgfmathprintnumber[fixed,precision=3]{\pgfplotspointmeta}},
nodes near coords style={font=\small},grid=major]
\addplot[fill=blue!35] coordinates {(0,0.014)(1,-0.580)(2,-0.699)};
\draw[red,dashed,thick] (axis cs:-0.45,0)--(axis cs:2.45,0);
\end{axis}\end{tikzpicture}\caption{240 MHz 最终时序闭合点与 300 MHz 两代实现的可用 WNS}\end{figure}


\PlaceholderFigure{250 MHz 时序报告截图}{请从 baseline-fabb599-clean250 归档截取 post-route physopt Timing Summary。}
\PlaceholderFigure{300 MHz 时序报告截图}{请从 backend-partition-fabb599 归档截取 routed 与 post-route physopt Timing Summary。}

\subsection{工程可复现性}

实现脚本统一完成 RTL 打包、Vivado 工程源刷新、IP/OOC 生成、布局布线、报告提取和归档。流程显式记录
代码提交、COE 身份、请求/实际频率、运行命令与关键报告；顶层实现并发不低于 16，IP/OOC 并发不低于
4。对于内联仿真模型，仿真保留模型，综合与打包阶段排除模型并绑定同名 Vivado IP，避免仿真结构与
物理实现的时延配置不一致。
